\section{From a desired load to emitter commands}\label{sec:inverse}

\subsection{A linear wave response produces a quadratic load}
Fix a geometry, material state and carrier frequency for which the linear wave
problem is well posed. Let $N$ be the number of independent complex mechanical
actuator channels. In one useful boundary representation, let $b_j(\bm x)$ be
the dimensionless normal-velocity pattern of channel $j$ on $\Gamma_a$.
The imposed peak velocity is $\sum_{j=1}^N b_jw_j$. The entries $w_j$ of
$\bm w\in\C^N$ have units of velocity and describe physical motion, not
uncalibrated electrical voltages.

Let $\mathcal A_\Gamma(\omega)$ be the linear wave operator, including all
interface and exterior conditions, and $\mathcal B_\Gamma(\omega)$ the source
operator. Its acoustic state $\bm z_a$ satisfies
\begin{equation}
 \mathcal A_\Gamma(\omega)\bm z_a
       =\mathcal B_\Gamma(\omega)\bm w.
 \label{eq:wave-operator}
\end{equation}
For each phase define a row response $\bm h_{p\ell}(\bm x)$ and a three-row
velocity response $\bm H_{v\ell}(\bm x)$ by
$P_\ell=\bm h_{p\ell}\bm w$ and
$\bm V_\ell=\bm H_{v\ell}\bm w$. The former has units
$\mathrm{Pa\,s\,m^{-1}}$ and the latter is dimensionless. Substitution in
\cref{eq:radiation-flux,eq:acoustic-traction} gives Hermitian matrices
\begin{align}
 \bm K_\ell(\bm x)&=
 \frac{\bm h_{p\ell}^\dagger\bm h_{p\ell}}{4\rho_\ell c_\ell^2}
 -\frac{\rho_\ell}{4}\bm H_{v\ell}^\dagger\bm H_{v\ell}
 +\frac{\rho_\ell}{2}\bm H_{v\ell}^\dagger\nn\nn^T\bm H_{v\ell},\\
 \bm K(\bm x;\Gamma,\omega)&=\bm K_- -\bm K_+,\qquad
 \boxed{\Pi(\bm x)=\bm w^\dagger\bm K(\bm x;\Gamma,\omega)\bm w.}
 \label{eq:quadratic-load}
\end{align}
A Hermitian matrix equals its conjugate transpose, so the load is real.
The matrix $\bm K$ has units $\mathrm{kg\,m^{-3}}$ and need not be positive
semidefinite for two interacting fluids. In the particular limit
\cref{eq:pressure-release-limit}, it is positive semidefinite at each point.
Tangential traction components also have Hermitian quadratic forms.

For a quiescent target, choose $M$ real, dimensionless virtual-displacement
functions $\phi_i$ compatible with the support and satisfying
$\int_\Gamma\phi_i\dd A=0$. Define generalized forces $F_i$, required forces
$b_i$, and integrated matrices $\bm K_i$ by
\begin{equation}
 \bm K_i=\int_\Gamma\phi_i\bm K\dd A,\qquad
 F_i=\bm w^\dagger\bm K_i\bm w,\qquad
 b_i=\int_\Gamma\phi_i(\sigma\kappa+\Delta\rho\Phi)\dd A.
 \label{eq:modal-loads}
\end{equation}
Here $\bm K_i$ has units $\mathrm{kg\,m^{-1}}$, while $F_i,b_i$ are forces.
The target equations are $F_i=b_i$ for all admitted variations. A finite
collection is a necessary projection of the continuum balance; its unsampled
complement and tangential balance still need verification. This construction
is reusable for different targets. The wave response must, however, be updated
when the target changes the scattering geometry.

\subsection{Feasibility, coherence and a useful relaxation}
Define the Hermitian drive matrix $\bm W=\bm w\bm w^\dagger$, whose entries
have units $\mathrm{m^2\,s^{-2}}$. A matrix $\bm W\succeq0$ is positive
semidefinite if $\bm q^\dagger\bm W\bm q\geq0$ for every complex vector
$\bm q$. A nonzero deterministic coherent drive has $\rank\bm W=1$.
Writing $\tr$ for trace, the exact fixed-target problem includes
\begin{equation}
 \tr(\bm K_i\bm W)=b_i,\qquad
 \bm W\succeq0,\quad\rank\bm W\leq1,\qquad
 W_{jj}\leq w_{j,\max}^2.
 \label{eq:lifted-feasibility}
\end{equation}
The positive bounds $w_{j,\max}$ specify peak channel velocities. The rank-zero
case is the undriven state. Let $F_{i,\rm ref}>0$ be meaningful force-error
scales, let $P_{\rm ref}>0$ be a power scale, and let $\beta_P\geq0$ weight
power. If the passive, fixed-geometry actuator model gives average input power
$P_{\rm in}=\bm w^\dagger\bm R_P\bm w$, the Hermitian matrix
$\bm R_P\succeq0$ has units $\mathrm{kg\,s^{-1}}$ and includes the phasor
averaging factor. A practical mismatch objective is
\begin{equation}
 \min_{\bm W}\ \frac12\sum_{i=1}^M
  \left[\frac{\tr(\bm K_i\bm W)-b_i}{F_{i,\rm ref}}\right]^2
       +\beta_P\frac{\tr(\bm R_P\bm W)}{P_{\rm ref}},
 \label{eq:sdr-cost}
\end{equation}
subject to the positivity, rank and hardware constraints. For example, a
pointwise peak pressure bound is linear in $\bm W$:
$\tr(\bm h_p^\dagger\bm h_p\bm W)\leq P_{\max}^2$, where $P_{\max}$ is a
declared pressure limit. Such bounds must hold throughout the specified domain,
not merely at convenient sample points.

Dropping the rank constraint gives a convex semidefinite relaxation when all
other constraints are convex at this fixed geometry. This is a standard
quadratic-optimization construction \cite{luo2010}. Its objective value is a
lower bound for the coherent problem, not proof that a physical coherent
solution attains it. If an optimum has rank one, its factor supplies the drive
up to a global phase. Explicitly, let $\lambda_1>0$ be the single nonzero
eigenvalue and $\bm q_1$ a unit eigenvector of a rank-one optimum. Let $A_j$
and $\varphi_j$ denote channel $j$'s peak velocity and phase, and $v_{a,j}(t)$
its real velocity coefficient. They follow as
\begin{equation}
 \bm w=\sqrt{\lambda_1}\bm q_1,\qquad
 A_j=|w_j|,\quad \varphi_j=\arg w_j,\qquad
 v_{a,j}(t)=A_j\cos(\omega t-\varphi_j).
 \label{eq:drive-recovery}
\end{equation}
The phase of a zero amplitude is immaterial. A higher-rank optimum requires
additional work.

In particular, writing $\bm W$ as a sum of rank-one factors does not create
independent same-frequency coherent fields: those fields add before their
intensity is squared. A higher-rank $\bm W$ can represent a physical covariance
$\mathbb E[\bm w\bm w^\dagger]$ only under an implemented incoherence or
time-multiplexing scheme with appropriate averaging, bandwidth and peak limits.
Covariance bounds constrain second moments; peak limits must additionally hold
for each physically realized waveform.
Switching among patterns must respect the acoustic and mechanical memory in
\cref{sec:acoustics}. Different carrier frequencies generally have different
$\bm K_i(\omega)$ and cannot be substituted into a fixed-frequency factorization
without changing the problem.

\begin{proposition}[A necessary load-feasibility certificate]
Let $\alpha_i$ be real coefficients and suppose
$\sum_i\alpha_i\bm K_i\succeq0$. Every exact feasible load vector in
\cref{eq:lifted-feasibility} must satisfy $\sum_i\alpha_i b_i\geq0$.
\end{proposition}
\begin{proof}
For $\bm W\succeq0$,
$\sum_i\alpha_i b_i=\tr[(\sum_i\alpha_i\bm K_i)\bm W]\geq0$.
\end{proof}
A violated inequality rules out even the relaxed load under that wave model.
The condition is only necessary; actuator bounds, coherence, optical
compatibility and stability can still rule out a load that passes it. Its
purpose is to expose unattainability before treating a failed optimizer as an
algorithmic problem.

\subsection{What changes when the surface is also unknown}
Let $\bm y$ collect the surface, hydrodynamic state, acoustic state and, where
needed, temperature and actuator variables. Let
$\mathcal R(\bm y,\bm w)=0$ collect the stationary governing equations,
volume conditions, boundary conditions and chosen pressure and phase gauges.
Choose nonnegative weights $\beta_W,\beta_d,\beta_s$ and a dimensionless
geometric mismatch $J_s$. A coupled design may minimize
\begin{equation}
 J_{\rm stat}=\beta_WJ_W+\beta_dJ_d+\beta_sJ_s
       +\beta_P P_{\rm in}/P_{\rm ref}.
 \label{eq:coupled-static-design}
\end{equation}
This minimization is subject to $\mathcal R=0$ and the stated optical,
actuator and stability constraints.
For a prescribed exact Cartesian surface, impose that target directly instead
of trading it away against an arbitrary geometric penalty. A tolerance-based
design must state its allowed optical error. The full problem is nonlinear
because both required curvature and acoustic transfer depend on geometry;
the fixed-target relaxation does not make it globally convex.

Let $\delta$ denote a first variation of geometry, material state or drive.
Differentiating \cref{eq:wave-operator} gives
\begin{equation}
 \mathcal A_\Gamma\,\delta\bm z_a
   =(\delta\mathcal B_\Gamma)\bm w
       +\mathcal B_\Gamma\delta\bm w
       -(\delta\mathcal A_\Gamma)\bm z_a.
 \label{eq:wave-sensitivity}
\end{equation}
Moving normals, surface measure, interface conditions and material coefficients
are part of this derivative. Freezing the acoustic field while differentiating
only capillary curvature omits the central wave--shape feedback.

If the linearized constrained stationary operator $\mathcal R_{\bm y}$ is
invertible on the gauge-fixed state space, implicit differentiation gives
\begin{equation}
 \delta\bm y=-\mathcal R_{\bm y}^{-1}
                         \mathcal R_{\bm w}\,\delta\bm w.
 \label{eq:implicit-sensitivity}
\end{equation}
At a fold or singular branch that inverse fails, and a local sensitivity cannot
be interpreted as a global design map. Resonance can also make it very poorly
conditioned.

Finally, if the commanded electrical variable is $\bm a$ and its loaded
mechanical response is $\bm w=\mathcal H(\Gamma,\omega)\bm a$, then holding
$\bm a$ fixed generally does not hold $\bm w$ fixed. The shape derivative
includes $\delta\mathcal H$. A stabilizing prediction for prescribed wall
velocity only applies to an experiment that actually controls that variable
or models the transducer dynamics needed to realize it.
